%======================================================================
\part*{Part III — Empirical Evaluation of IHS for Weighted CSPs (Petrova, Larrosa, Roll\'on, 2025)}
\addcontentsline{toc}{part}{Part III — Empirical Evaluation of IHS for Weighted CSPs}
%======================================================================

\begin{abstract}
This part summarizes the method and findings of~\citet{petrova2025empirical},
who carry out an empirical study of IHS for the Weighted CSP framework. Where
Parts~I--II frame IHS theoretically as Dantzig--Wolfe on the dual packing LP,
Petrova et~al.\ vary the two main algorithmic axes (computation of hitting
vectors and core improvement) and report on 32 configurations across a
heterogeneous benchmark suite. The two studies are complementary: the OR
framing in Parts~I--II predicts that techniques relaxing the hitting-vector
subproblem and that strengthen cores should both pay off; the empirical study
confirms this is benchmark-dependent but identifies cost-function merging
combined with maximal-core extraction as the most robust default.
\end{abstract}

%----------------------------------------------------------------------
\section{Problem setting: Weighted CSP}
\label{sec:wcsp}
%----------------------------------------------------------------------

A \emph{Weighted Constraint Satisfaction Problem} (WCSP) is a triple
$(X, C, F)$ where $X$ is a set of variables over finite domains, $C$ is a set
of hard constraints (Boolean functions over scopes of $X$), and
$F = \{f_1, \ldots, f_m\}$ is a set of \emph{cost functions}, each mapping
assignments of its scope to a nonnegative cost. A solution is a complete
assignment satisfying all constraints; its cost is the sum of cost-function
values; the goal is a minimum-cost solution $w^\star$. WCSP generalizes
weighted MaxSAT: each soft clause becomes a unary cost function over its
falsification indicator, and hard clauses move into $C$.

%----------------------------------------------------------------------
\section{Vector-valued IHS for WCSP}
\label{sec:vecihs}
%----------------------------------------------------------------------

To handle multi-valued cost functions cleanly the paper recasts IHS in
vector-valued language. A \emph{cost vector}
$\vec{v} = (v_1, \ldots, v_m)$ assigns to each cost function $f_i$ a value
$v_i$ that must be an actual cost occurring in $f_i$. The vector induces a
CSP $(X, C \cup F_{\vec{v}})$ in which each soft cost function $f_i$ is
replaced by the constraint $f_i \le v_i$. If this CSP is satisfiable, $\vec{v}$
is a \emph{solution vector}; otherwise it is a \emph{core}. Costs and
dominance lift componentwise: $\vec{u} \le \vec{v}$ holds iff $u_i \le v_i$ for
all $i$; a set of vectors $\mathcal{V}$ \emph{dominates} a vector $\vec{u}$
if some $\vec{v} \in \mathcal{V}$ has $\vec{v} \le \vec{u}$. A vector that is
not dominated by a set of cores $\mathcal{K}$ is said to \emph{hit}
$\mathcal{K}$; the \emph{minimum-cost hitting vector} (MHV) of $\mathcal{K}$ is
a hitting vector of minimum cost. Computing MHV is NP-hard (it reduces to
minimum hitting set), so it is the analog of the IHS master.

Two structural bounds drive the algorithm. For any solution vector $\vec{h}$
and any set of cores $\mathcal{K} \subseteq \mathit{Cores}$,
\[
\mathit{MHV}(\mathcal{K}) \;\le\; w^\star \;\le\; \mathit{cost}(\vec{h}).
\]
The algorithm terminates when these two bounds meet. Termination is
characterized in terms of \emph{goal cores}: $\mathit{MHV}(\mathcal{K}) = w^\star$
iff $\mathcal{K}$ dominates every \emph{maximal core of cost less than
$w^\star$}. The IHS task is therefore to compute a set $\mathcal{K}$ that
dominates every such goal core and an optimal solution.

%----------------------------------------------------------------------
\section{Baseline algorithm}
\label{sec:baseline}
%----------------------------------------------------------------------

The baseline IHS algorithm (\textsc{IHS-lb}) is given in
Algorithm~\ref{alg:ihs-lb}. It maintains a working core set $\mathcal{K}$ with
a lower bound $\mathit{lb}$ and an upper bound $\mathit{ub}$. At each
iteration it computes an MHV $\vec{h}$ of $\mathcal{K}$ (raising
$\mathit{lb}$), solves the induced CSP and, if unsatisfiable, calls a routine
$\textsc{ImprCore}$ that grows $\vec{h}$ into an improved core
$\vec{k} \ge \vec{h}$, which is added to $\mathcal{K}$.

\begin{algorithm}[t]
\caption{Baseline IHS for WCSP (\textsc{IHS-lb}). $\textsc{MinCostHV}(\mathcal{K})$ returns a minimum-cost hitting vector; $\textsc{ImprCore}$ improves a hitting vector that is a core into a larger core.}
\label{alg:ihs-lb}
\begin{algorithmic}[1]
\Require WCSP $(X, C, F)$
\Ensure cost $w^\star$ of an optimal solution vector
\State $\mathcal{K} \gets \emptyset$;\ $\mathit{lb} \gets 0$;\ $\mathit{ub} \gets \infty$
\While{$\mathit{lb} < \mathit{ub}$}
  \State $\vec{h} \gets \textsc{MinCostHV}(\mathcal{K})$;\quad $\mathit{lb} \gets \mathit{cost}(\vec{h})$
  \If{$\textsc{SolveCSP}(X,\, C \cup F_{\vec{h}})$}
      \State $\mathit{ub} \gets \mathit{cost}(\vec{h})$
  \Else
      \State $\vec{k} \gets \textsc{ImprCore}(X, C, F, \mathit{ub}, \vec{h})$;\quad $\mathcal{K} \gets \mathcal{K} \cup \{\vec{k}\}$
  \EndIf
\EndWhile
\State \Return $\mathit{lb}$
\end{algorithmic}
\end{algorithm}

A symmetric variant \textsc{IHS-ub} computes a \emph{cost-bounded} hitting
vector (any hitting vector of cost less than the current $\mathit{ub}$,
solved as a decision problem) instead of the cost-minimal one, prioritizing
upper-bound tightening over lower-bound tightening. Each iteration of either
variant decomposes into two NP-hard pieces: the hitting-vector subproblem
(integer optimization) and the induced CSP (decision).

%----------------------------------------------------------------------
\section{Algorithmic alternatives (the 32 implementations)}
\label{sec:alts}
%----------------------------------------------------------------------

The paper isolates two design axes and crosses them with one preprocessing
choice, yielding $4 \times 4 \times 2 = 32$ implementations.

\subsection{Axis 1 --- Computation of hitting vectors (4 options)}

\begin{itemize}
\item \textbf{\textsc{MinCostHV} (IHS-\emph{lb}).} Optimal hitting vector via
0/1 integer programming. Each iteration strictly tightens the lower bound;
expensive per iteration.

\item \textbf{\textsc{CostBoundedHV} (IHS-\emph{ub}).} Any hitting vector of
cost less than the current $\mathit{ub}$. A decision problem, hence cheaper;
but the cost is not guaranteed to be below $w^\star$, so new cores need not
contribute to the lower bound (cf.\ the goal-cores characterization above).

\item \textbf{Greedy variant in \emph{lb} mode (IHS-\emph{grdlb}).} An
inexact, heuristic hitting-vector computation that starts from
$\vec{h} = \vec{0}$ and increments components by a ratio criterion (analogous
to greedy vertex cover). When a greedy step fails to make progress the next
iteration falls back to \textsc{MinCostHV}, preserving correctness.

\item \textbf{Greedy variant in \emph{ub} mode (IHS-\emph{grdub}).} As above,
but with fallback to \textsc{CostBoundedHV}.
\end{itemize}

The intent of relaxing \textsc{MinCostHV} to a cost-bounded or greedy variant
is the same that motivates in-out separation and dual stabilization in
column generation (Part~I, \S\ref{sec:techniques}): cheaper, less-optimal
probe points trade per-iteration cost for iteration count.

\subsection{Axis 2 --- Core improvement (\textsc{ImprCore}, 4 options)}

All four variants take a hitting vector $\vec{h}$ that is a core and grow it
into a larger core $\vec{k} \ge \vec{h}$ by incrementing components while
preserving the core condition.

\begin{itemize}
\item \textbf{Maximal cores.} Strongest and most expensive: keep incrementing
until no further component can be increased without breaking the core
condition. Analogous to destructive MUS extraction. Produces cores that
dominate the largest set of goal cores per call.

\item \textbf{Lazy cores.} Weakest and cheapest: no explicit improvement. The
assumption-based SAT solver already returns a core that is implicitly
improved at no extra cost.

\item \textbf{Cost-bounded cores.} Intermediate: improve $\vec{h}$ until its
cost reaches $\mathit{ub}$. If $\mathit{ub} > w^\star$, the resulting core
contributes to dominating goal cores; otherwise the improvement is wasted.

\item \textbf{Partially maximal cores.} Soften maximality to require only that
the increment of \emph{some} single component produces a solution vector
(rather than every component). Used in earlier work~\citep{moreno2013implicit}.
\end{itemize}

Across all four variants the implementation prefers, at each candidate step,
the component with the lowest value~$v_i$: having $\mathcal{K}$ cores whose
minimum value is as large as possible makes them more costly to hit,
accelerating the lower bound.

The maximal-cores option is the analog of Pareto-optimal core selection in
Part~I (\S\ref{sec:techniques}, \citet{magnanti1981accelerating}): both
strengthen the cuts/columns generated per oracle call at a higher
per-iteration price.

\subsection{Axis 3 --- Cost-function merging (with vs.\ without)}

A preprocessing step (see~\citet{petrova2025empirical} and references therein)
heuristically clusters cost functions along a tree decomposition and
``virtually'' merges each cluster into a single cost function with extended
domain. This shrinks the dimensionality of the cost-vector space and is shown
empirically to be the single most impactful design choice.

\subsection{Implementation choices}

\begin{itemize}
\item \textsc{MinCostHV} and \textsc{CostBoundedHV} are modelled as 0/1 IPs
and solved with \emph{CPLEX}.
\item Induced CSPs are encoded as CNF and solved with the assumption-based
SAT solver \emph{CaDiCaL}; this is what makes ``lazy cores'' free.
\item All instances are pre-processed and made \emph{virtually arc consistent}
(VAC) before solving.
\item Some iterations are allowed to extract \emph{disjunctive cores} (cores
covering multiple branches simultaneously), as in~\citet{allouche2015anytime}.
\item Time limit: 1 hour per instance.
\end{itemize}

%----------------------------------------------------------------------
\section{Empirical landscape (high-level)}
\label{sec:wcspresults}
%----------------------------------------------------------------------

The benchmark suite (12 classes) spans uniform random, scale-free random, and
structured problems (EHI, SPOT5, driverlog, Grid, Normalized, Pedigree,
CELAR), with up to a few thousand variables and cost functions per instance.

\paragraph{Cost-function merging is dominant.} Speed-ups range from $1.3\times$
(Normalized) to $268\times$ (random domains) with only the Grid class harmed
by merging---suggesting the tree-decomposition heuristic is unsuited to
its regular structure.

\paragraph{No algorithm uniformly dominates.} Different benchmark classes
prefer dramatically different combinations, with ratios up to several orders
of magnitude between the best and the worst configuration on the same
benchmark. The authors warn against concluding that IHS ``is not suitable''
for any particular benchmark class without trying multiple combinations.

\paragraph{Maximal cores are the most common best option.} For $10$ of $12$
benchmark classes maximal cores produce the best running time. Pedigree and
random-domains are the exceptions: there the induced CSPs are SAT-hard
enough that the cheaper lazy cores pay off.

\paragraph{The intermediate options are never the best for core improvement.}
Cost-bounded cores turn out to be empirically close to lazy cores (the
assumption-based SAT solver already provides cores whose cost reaches
$\mathit{ub}$); partially maximal cores tend to be too weak. Only the two
extremes (maximal or lazy) ever win.

\paragraph{Greedy hitting vectors rarely pay off.} Except in the very early
iterations, greedy variants produce cheap-but-useless iterations and end up
converging to their exact counterpart; better heuristics could plausibly
change this picture.

\paragraph{Competitiveness with the state-of-the-art.} Compared with
Toulbar2 (the leading WCSP solver), IHS still trails on most random
benchmarks but is competitive---and outperforms Toulbar2---on benchmarks with
small domains and few distinct costs (EHI, SPOT5, Grid, Normalized,
Pedigree). These two structural features (small domains, few distinct
costs) appear to be the most reliable proxy for IHS suitability.

%----------------------------------------------------------------------
\section{Connection to Parts~I and~II}
\label{sec:wcspconnect}
%----------------------------------------------------------------------

The two design axes Petrova et~al.\ vary line up cleanly with the OR-derived
techniques framed in Part~I (\S\ref{sec:techniques}).

\begin{itemize}
\item \textbf{Axis 1 (hitting-vector relaxation) $\leftrightarrow$
in-out separation / cost-bounded probing.} The paper's \textsc{CostBoundedHV}
and greedy options are concrete realizations of the idea that the master
need not be solved to optimality at every iteration; the same observation
underlies in-out separation in column generation~\citep{benameur2007stabilization}.

\item \textbf{Axis 2 (core improvement) $\leftrightarrow$ Pareto-optimal
core selection.} The maximal-cores option of Petrova et~al.\ is the
non-fractional analog of the Magnanti--Wong dominance criterion: both
strengthen the cuts/columns produced per oracle call.

\item \textbf{Axis 3 (cost-function merging) $\leftrightarrow$ Dantzig--Wolfe
preprocessing.} Tree-decomposition-based merging compresses the
constraint structure before the IHS loop, in the spirit of OR
preprocessing techniques (e.g.\ aggregation in column generation) that
reduce master dimension prior to solving.
\end{itemize}

Two features of the paper depart from the Parts~I--II framing. First, the
paper does not consider an \emph{LP-relaxation} view of the hitting-vector
subproblem, hence none of its variants exploit fractional dual prices
$\bar{y}$; the price-aware extraction and Wentges stabilization techniques
of Part~I are therefore complementary rather than overlapping. Second, the
paper does \emph{not} consider warm-starting from a complementary algorithm
(e.g.\ Toulbar2's incremental lower bound, or a short core-guided phase);
cross-family warm-starting (Part~I, \S\ref{sec:techniques}) is an
unexplored direction.

The empirical evidence on the cost-function merging axis is also a
cautionary note for theoretical analyses that treat the cost-function
collection as fixed: merging changes the effective dimensionality of the
hitting-vector space, which Parts~I--II do not model and which dominates
solver performance.

\paragraph{Summary.} Parts~I--II provide a structural reason to expect that
relaxed hitting-vector subproblems and stronger cores both help; the empirical
study of~\citet{petrova2025empirical} confirms the latter (maximal cores) and
qualifies the former (cost-bounded helps mainly via assumption-based SAT,
greedy heuristics need better design). The complementarity is clean: the OR
framing identifies the design axes worth varying, and the empirical study
populates them with measured trade-offs.


